\clearpage\fancyhead[L]{\small VESTA.}
\begin{center}
VESTA.
\end{center}

\begin{center}
\rule{0.6\linewidth}{0.4pt}\\
Monatliche Correspondenz. Vol. XV. Pp. 590–600. June 1807.\\
\rule{0.6\linewidth}{0.4pt}
\end{center}

--- --- ---

\par\indent Dr. \textsc{Gauss} in Brunswick made the following two observations, based on numerous comparisons with \textsc{Piazzi}’s star, R.A. $180^\circ\,51'\,27''5$, northern declination $12^\circ\,35'\,14''0$.

\begin{center}\begingroup
\fontsize{10}{12}\selectfont
\setlength{\tabcolsep}{2.5pt}\renewcommand{\arraystretch}{1.15}
\begin{tabular}{c|c|c|c}
1807. & \shortstack{Mean time in\\Brunswick} & \shortstack{Apparent right\\ascension \VestaRATwoNinety{}} & \shortstack{Apparent northern\\declination \VestaDecTwoNinety{}} \\ \hline
April 6 & $8^{\text{u}}\,22^{\text{m}}\,45^{\text{s}}$ & $182^\circ\,21'\,54''6$ & $12^\circ\,28'\,0''9$ \\
8 & $8\ \ 28\ \ 49$ & $181\ \ 56\ \ 51.0$ & $12\ \ 35\ \ 45.8$
\end{tabular}\endgroup\end{center}

\par\indent One would have expected the admirable Dr. \textsc{Gauss}, who has earned such unrivaled distinction for his work on all the newly discovered small planets, to retain that distinction with \emph{Vesta}. This time, however, the great geometer has surpassed himself. “I have often admired our excellent \textsc{Gauss},” (Dr. \textsc{Olbers} writes to us on 20 April) “but this time he has astonished me. He received my observation of \emph{Vesta} from 17 April at $9\tfrac{1}{2}$ in the evening on 20 April, and by 5 in the afternoon on 21 April he could already send me the calculated elliptic elements and the comparison with all the observations. For all this-
\SourceLeaf{291}%
he needed only ten hours. This is also proof of the great flexibility of his new and greatly improved method\HistoricalNote{27} of determining planetary orbits.”

\par\indent Dr. \textsc{Gauss} himself writes to us about this. “I undertook the first determination of the orbit on 20 April, when I had received Dr. \textsc{Olbers}’s observation of 17 April, which brought the interval to 19 days. I sent the results to Dr. \textsc{Olbers} on 21 April. But the interval was still so short, almost all the observations had been made with only the ring micrometer, and the small inclination of this planet’s orbit also makes the results relatively less reliable. I did not yet dare to present them as a genuine approximation to the truth, so I held them back with the intention of improving them once more using later observations. Only yesterday (27 April) did more recent observations reach me, and I thought I would use them at once for a new calculation. To my great pleasure, however, I find that my first results still agree with them so closely that I really would not know what to change with any confidence yet. So I am postponing the new calculation until the error has grown to a distinctly noticeable size. I will wait no longer to send you the I. elements of \emph{Vesta}.

\begin{center}\begingroup\fontsize{9}{11}\selectfont
\begin{tabular}{ll}
\hypertarget{vesta1807-first-elements}{}\shortstack[l]{Epoch of longitude, 29 March 1807\\$12^{\text{u}}$ M.\,T. in Bremen} . & $193^\circ\,8'\,4''6$ \\
Aphelion .\ .\ .\ .\ .\ .\ .\ .\ .\ .\ .\ .\ .\ .\ .\ . & $69^\circ\;7\;40.9$ \\
Ascending node .\ .\ .\ .\ .\ .\ .\ .\ .\ .\ .\ .\ . & $103\;8\;36.2$
\end{tabular}\endgroup\end{center}

\noindent All three are measured from the mean equinox, and the last two are treated as fixed in tropical coordinates.

\begin{center}\begingroup\fontsize{9}{11}\selectfont
\begin{tabular}{ll}
Mean daily motion .\ .\ .\ .\ .\ .\ .\ .\ .\ .\ .\ . & $978''909$ \\
Eccentricity .\ .\ .\ .\ .\ .\ .\ .\ .\ .\ .\ .\ .\ .\ .\ .\ .\ . & $0.097505$ \\
Logarithm of the semimajor axis .\ .\ .\ .\ .\ .\ .\ .\ . & $0.3728428$ \\
Inclination of the orbit .\ .\ .\ .\ .\ .\ .\ .\ .\ .\ .\ .\ .\ .\ . & $7^\circ\;5'\;49''5$
\end{tabular}\endgroup\end{center}

\par\indent “Here is the comparison with all the observations that have become known.

\begin{center}\begingroup
\fontsize{8}{10}\selectfont
\setlength{\tabcolsep}{2.5pt}\renewcommand{\arraystretch}{1.15}
\begin{tabular}{c|c|c|c|c|c|c}
1807 & \shortstack{Mean time\\at observing\\site} & \shortstack{Calculated right\\ascension\\\VestaRATwoNinetyOne{}} & \shortstack{Difference.} & \shortstack{Calculated\\northern declination\\\VestaDecTwoNinetyOne{}} & \shortstack{Difference.} & Observer. \\ \hline
March 29 & $8^{\text{u}}\,21^{\text{m}}\,20^{\text{s}}$ & $184^\circ\,8'\,54''4$ & $-1''4$ & $.\ .\ .\ .\ .$ & $.\ .\ .$ & Olbers \\
        & $10\ \ 31\ \ 16$ & $184\ \ 7\ \ 40.2$ & $-10.6$ & $11^\circ\,47'\,52''4$ & $+5''4$ & \\
30 & $8\ \ 34\ \ 53$ & $183\ \ 55\ \ \ 3.0$ & $+6.2$ & $.\ .\ .\ .\ .$ & $.\ .\ .$ & \\
   & $8\ \ 44\ \ \ 8$ & $183\ \ 54\ \ 57.7$ & $+11.9$ & $11\ \ 53\ \ 18.6$ & $+7.6$ & \\
   & $12\ \ 33\ \ 17$ & $183\ \ 52\ \ 44.2$ & $+3.4$ & $11\ \ 54\ \ 14.7$ & $-12.3$ & \\
April 1 & $9\ \ 50\ \ \ 0$ & $183\ \ 27\ \ \ 3.7$ & $+0.9$ & $12\ \ \ 4\ \ 40.7$ & $-11.3$ & \\
   & $12\ \ \ 4\ \ 35$ & $183\ \ 25\ \ 47.8$ & $+23.6$ & $12\ \ \ 5\ \ \ 9.6$ & $-22.1$ & Bessel \\
2 & $8\ \ 21\ \ \ 1$ & $183\ \ 14\ \ 23.8$ & $+6.0$ & $12\ \ \ 9\ \ 34.7$ & $-12.3$ & Olbers \\
3 & $8\ \ 16\ \ 49$ & $183\ \ \ 1\ \ \ 3.6$ & $+2.8$ & $12\ \ 14\ \ 33.6$ & $-6.4$ & \\
4 & $9\ \ \ 9\ \ 43$ & $182\ \ 47\ \ 20.9$ & $+0.1$ & $12\ \ 19\ \ 29.1$ & $-15.9$ & \\
5 & $11\ \ 17\ \ \ 3$ & $182\ \ 33\ \ 10.7$ & $-0.2$ & $12\ \ 24\ \ 21.6$ & $+2.5$ & Harding \\
6 & $8\ \ 22\ \ 45$ & $182\ \ 21\ \ 48.5$ & $-6.1$ & $12\ \ 28\ \ \ 6.6$ & $+5.7$ & Gauss \\
  & $11\ \ 12\ \ 15$ & $182\ \ 20\ \ 21.6$ & $-12.2$ & $12\ \ 28\ \ 34.6$ & $+45.2$ & Harding \\
8 & $8\ \ 21\ \ 37$ & $181\ \ 56\ \ 42.9$ & $-0.1$ & $12\ \ 35\ \ 50.1$ & $-8.0$ & Olbers \\
  & $8\ \ 28\ \ 49$ & $181\ \ 56\ \ 42.2$ & $-8.8$ & $12\ \ 35\ \ 50.1$ & $+4.3$ & Gauss \\
  & $11\ \ \ 2\ \ 43$ & $181\ \ 55\ \ 22.5$ & $+3.3$ & $12\ \ 36\ \ 13.6$ & $+1.4$ & Harding
\end{tabular}\endgroup\end{center}
\newpage
\begin{center}\begingroup
\fontsize{8}{10}\selectfont
\setlength{\tabcolsep}{2.5pt}\renewcommand{\arraystretch}{1.15}
\begin{tabular}{c|c|c|c|c|c|c}
\SourceLeaf{292}%
1807 & \shortstack{Mean time\\at observing\\site} & \shortstack{Calculated right\\ascension\\\VestaRATwoNinetyTwo{}} & \shortstack{Difference.} & \shortstack{Calculated\\northern declination\\\VestaDecTwoNinetyTwo{}} & \shortstack{Difference.} & Observer. \\ \hline
April 9 & $8^{\text{u}}\,30^{\text{m}}\,46^{\text{s}}$ & $181^\circ\,44'\,25''6$ & $+8''6$ & $12^\circ\,39'\,19''8$ & $+5''7$ & Olbers \\
   & $10\ \ 57\ \ 58$ & $181\ \ 43\ \ 13.8$ & $+11.5$ & $12\ \ 39\ \ 39.6$ & $+3.5$ & Harding \\
12 & $8\ \ 27\ \ 20$ & $181\ \ \ 9\ \ 25.6$ & $-8.4$ & $12\ \ 48\ \ \ 7.4$ & $-8.7$ & Olbers \\
13 & $10\ \ 39\ \ 14$ & $180\ \ 57\ \ 28.3$ & $+11.3$ & $12\ \ 50\ \ 41.5$ & $-1.5$ & Bode \\
14 & $8\ \ 19\ \ 51$ & $180\ \ 47\ \ 36.3$ & $-21.7$ & $12\ \ 52\ \ 37.2$ & $-10.9$ & Olbers \\
   & $10\ \ 34\ \ 33$ & $180\ \ 46\ \ 39.2$ & $-0.7$ & $12\ \ 52\ \ 44.4$ & $+4.7$ & Harding \\
17 & $8\ \ 12\ \ 35$ & $180\ \ 17\ \ 21.6$ & $-5.2$ & $12\ \ 57\ \ 18.3$ & $+2.3$ & Olbers \\
23 & $8\ \ 30\ \ 52$ & $179\ \ 26\ \ 39.7$ & $-1.3$ & $12\ \ 59\ \ 20.4$ & $+0.4$ & \\
25 & $9\ \ 53\ \ 49$ & $179\ \ 13\ \ \ 0.6$ & $-18.4$ & $12\ \ 57\ \ 52.4$ & $-5.6$ & \\
   & $9\ \ 49\ \ 26$ & $179\ \ 12\ \ 57.1$ & $+8.1$ & $12\ \ 57\ \ 51.9$ & $-14.1$ &
\end{tabular}\endgroup\end{center}

\par\indent This agreement is as good as could be expected with these observations. As the interval now amounts to 27 days, Dr. \textsc{Gauss} thinks that these elements give at least some idea of the principal dimensions of the orbit. \emph{Vesta} would therefore complete its orbit much faster than the other small planets. Its orbital inclination would be the smallest of all and its eccentricity moderate. The latter circumstance and the planet’s great brightness give Dr. \textsc{Gauss} very strong hopes that we will find earlier observations of this new planet, perhaps even in \textsc{Flamsteed}. As soon as the elements are known with a little more certainty, Dr. \textsc{Gauss} will look into this.

\par\indent On 18 May this incomparable astronomer writes to us. “The first elements of \emph{Vesta} that I sent you some time ago, based on 19 days of observations, continued to agree very well with the observations until toward the end of April. In May, however, the difference began to become somewhat noticeable, and on 6 May it had grown to $+33''$ in right ascension and $-41''$ in declination. An improvement of the elements was hardly necessary yet. What chiefly prompted me was that I had begun to suspect certain stars in older catalogues and was impatient to get an idea of how large a change the elements might still require. On 9 May, after receiving Dr. \textsc{Olbers}’s observation of the $6^{\text{th}}$, I undertook a second calculation of the orbit, and I have the honor of reporting its results here.

\Needspace{16\baselineskip}
\begin{center}
II. \emph{Elements of Vesta}
\end{center}

\begin{center}
Based on 38 days of observations.
\end{center}

\begin{center}\begingroup\fontsize{9}{11}\selectfont
\begin{tabular}{ll}
\shortstack[l]{Epoch of longitude, 29 March 1807 $12^{\text{u}}$\\M.\,T. in Bremen} .\ .\ .\ .\ .\ . & $192^\circ\,9'\,53''9$ \\
Aphelion .\ .\ .\ .\ .\ .\ .\ .\ . & $69\;57\;52$ \\
Ascending node .\ .\ .\ .\ .\ . & $103\;18\;34$ \\
Mean daily tropical motion  & $980''707$ \\
Eccentricity .\ .\ .\ .\ .\ .\ .\ .\ . & $0.0872230$ \\
Logarithm of the mean distance . & $0.3723521$ \\
Inclination of the orbit .\ .\ .\ .\ .\ .\ . & $7^\circ\,8'\,6''8$
\end{tabular}\endgroup\end{center}

\par\indent “These elements still agree beautifully with the observations through the $11^{\text{th}}$, as far as the observations I have received extend. I have compared 68 observations in all so far. The agreement is so close throughout that nowhere can a definite difference be stated between the calculated position and the true position.”

\par\indent “I still have great hopes that we will find observations of \emph{Vesta} in the \emph{Historia coelestis}. For the missing stars that I marked as suspicious right at the start, however, that hope has mostly vanished. I would still reserve judgment only on 91 \emph{Virginis} as given by
\SourceLeaf{293}%
\textsc{Flamsteed}. The longitude would agree very well with the position of \emph{Vesta} if the mean daily motion were only about $6''$ greater than according to the II. elements, but the latitude \emph{seems} a little too large. We will soon be able to judge the matter more definitely.”

--- --- ---
\hfill [\textsc{von Zach}.]
